\section{Key Stakeholders and Bloc Positions}

\subsection{European Union and Western Allies}

This group believes in strong accountability measures and are entirely pro-regulation. The European bloc advocates for criminal prosecution and binding state obligations when cultural heritage is destroyed as a deliberate weapon of war.\footnote{Dijkstal, H. J. (2019). Destruction of cultural heritage before the ICC. Journal of International Criminal Justice, 17(2), 391–412. \url{https://doi.org/10.1093/jicj/mqz017}} This bloc has historically been involved or even the drivers of the development of the legal instruments protecting cultural heritage.\footnote{Boylan, P. J. (1993). Review of the Convention for the Protection of Cultural Property in the Event of Armed Conflict (The Hague Convention of 1954). UNESCO. \url{https://unesdoc.unesco.org/ark:/48223/pf0000102816}} The EU has even adopted a Council Resolution (2019/880) to declare how important cultural goods are and to prevent trafficking of items of cultural heritage. The bloc stands by the position that expanding ICC jurisdiction, mandatory reporting from market actors, and proactive digital surveillance of online antiquities sales is the solution.\footnote{European Parliament. (2023). Protecting cultural heritage from armed conflicts in Ukraine and beyond. Policy Department for Structural and Cohesion Policies. \url{https://www.europarl.europa.eu/RegData/etudes/STUD/2023/733120/IPOL_STU\%282023\%29733120_EN.pdf}} Accountability must be punitive and reconstructive for the community as a whole.\footnote{Wierczyńska, K., \& Jakubowski, A. (2017). Individual responsibility for deliberate destruction of cultural heritage: Contextualising the ICC judgment in the Al-Mahdi case. Chinese Journal of International Law, 16(4), 695–721. \url{https://doi.org/10.1093/chinesejil/jmx029}} The European bloc argues that the way forward is a robust international accountability framework, based on the “shared interest of humanity”.\footnote{Francioni, F. (2004). Beyond state sovereignty: The protection of cultural heritage as a shared interest of humanity. \emph{Michigan Journal of International Law}, 25(4), 1209–1228}

\subsection{United States and Common Law Allies}

This bloc shares the idea that condemnation is required for the destruction of cultural heritage, however, they diverge from the European bloc on the appropriate measures. They prioritise national sovereignty and voluntary cooperation, alongside market-based solutions.\footnote{Forrest, C. (2010). \emph{International law and the protection of cultural heritage}. Routledge.} They are sceptical of enforcement mechanisms which require international judicial proceedings. Some of these states have not ratified the Second Protocol to the 1954 Hague Convention due to this reason.\footnote{UNESCO. (2024). States parties to the Second Protocol. \url{https://en.unesco.org/protecting-heritage/convention-and-protocols/states-parties}} This bloc prefers to use sanctions and trade restrictions to combat the destruction of global heritage.\footnote{United Nations Security Council. (2015). Resolution 2199 (S/RES/2199). \url{https://undocs.org/S/RES/2199(2015)}} This bloc focuses on preventing and monitoring the artifact trafficking avenues and how they finance terrorist networks.\footnote{Brodie, N. (2015). Syria and its regional neighbours: A case of cultural property protection policy failure? International Journal of Cultural Property, 22(2–3), 317–335. \url{https://doi.org/10.1017/S0940739115000132}} This involves working with other enforcement agencies or committees such as INTERPOL or the United Nations Office on Drugs and Crime (UNODC).

\subsection{States Directly Affected by Heritage Destruction}

This bloc consists of states whose cultural heritage has been directly affected or damaged by armed conflict. These states are often the ones who advocate the most for international accountability, reconstruction, and recognition as a trifold framework.\footnote{International Criminal Court. (2017). Prosecutor v. Ahmad Al Faqi Al Mahdi, Reparations Order (ICC-01/12-01/15-236). \url{https://www.icc-cpi.int/mali/al-mahdi}} This bloc often supports language and statements which characterise the intentional destruction of cultural heritage as an attack on human dignity.\footnote{UN Special Rapporteur on Cultural Rights. (2016). Report on the intentional destruction of cultural heritage (A/71/317). United Nations. \url{https://undocs.org/A/71/317}} This bloc supports continuous funding mechanisms for occurrences of reconstruction.\footnote{Isakhan, B., \& Meskell, L. (2019). UNESCO's project to 'Revive the Spirit of Mosul': Iraqi and Syrian opinion on heritage reconstruction after the Islamic State. International Journal of Heritage Studies, 25(11), 1189–1204. \url{https://doi.org/10.1080/13527258.2019.1578988}} These states often require criminal prosecution, state responsibility for the failure to protect and reparations for affected communities to be part of the accepted solutions moving forward.

\subsection{States with Sovereignty Concerns}

This bloc is formed of states who endorse protection of cultural heritage but oppose any systematic accountability mechanism which could be used against state actors.\footnote{Schabas, W. A. (2016). The International Criminal Court: A commentary on the Rome Statute (2nd ed.). Oxford University Press.} Another mechanism these states oppose is establishing precedent in international courts that allow for scrutiny or examination of how states govern heritage issues within their own borders.\footnote{International Law Commission. (2001). Articles on Responsibility of States for Internationally Wrongful Acts. United Nations. \url{https://legal.un.org/ilc/texts/instruments/english/draft_articles/9_6_2001.pdf}} The core belief of this bloc is that cultural heritage should be protected, but not at the cost of state sovereignty. These states believe that each issue should be governed domestically, and do not want ICC jurisdiction to be expanded for risk of foreign intervention into state matters. These states also resist international verification mechanisms regarding documentation and monitoring destruction or damage to cultural heritage.\footnote{Lostal, M. (2017). International cultural heritage law in armed conflict: Case studies of Syria, Libya, Mali, the invasion of Iraq, and the Buddhas of Bamiyan. Cambridge University Press} This bloc believes the way to move forward is to use a framework of independent state implementation.

\subsection{Developing States}

This bloc views accountability while focusing on highlighting structural inequality. These states often argue that meaningful heritage protection requires them to address the disparities in institutional capacity, technologies and funding available to different states and organisations. Developing states often suffer disproportionately more from destruction of heritage, as they are more likely to have active armed conflicts, and they possess the fewest resources to combat the destruction as it occurs and afterwards. The destruction of these states' heritage is often undocumented and not prioritised for repair or discussion on rebuilding, so they often feel left out of the conversation. This bloc highly supports the guidelines by the World Customs Organisation in 2020 even as they require sustained investments into training and technology, and further institutional development.
